\subsection{Full-GT stability of the new Hurwitz representations}

\begin{statement}{D23.P18}{Proved}{Functorial stability of an isolated finite Hurwitz orbit}
\end{statement}
Let $G$ be a finite group, let $\mathcal C$ be a conjugacy class
in $G$, and let $\mathcal N_n=\mathcal C^n/G$ carry the left
Hurwitz action
\[
 (g_i,g_{i+1})\longmapsto(g_i g_{i+1}g_i^{-1},g_i).
\]
Every multiplier-one element $h\in\widehat{GT}$ induces a
permutation $V_h$ of $\mathcal N_n$ satisfying
\[
 V_h q_n(\beta)V_h^{-1}=q_n(\Phi_h(\beta))
 \qquad(\beta\in\widehat B_n),
 \tag{D23.P18.1}
\]
where $\Phi_h$ is the standard Drinfeld automorphism. Therefore
an orbit whose degree occurs exactly once among all
$B_n$-orbits in $\mathcal N_n$ is preserved. Its permutation
representation has kernel invariant under the multiplier-one
subgroup of full GT, and its transform
under $\Phi_h$ is conjugate to the original representation by
an actual permutation of that orbit.

In particular, the degree-1344 representation of D23.P17 is
stable in this sense. If $F=f(X_1,X_2)$ and $R=f(X_3,X_4)$ in
its permutation image, there exists $V\in S_{1344}$ with
\[
 \boxed{\quad
 VAV^{-1}=A,\qquad
 VBV^{-1}=F^{-1}BF,\qquad
 VCV^{-1}=RCR^{-1}.
 \quad}                                                     \tag{D23.P18.2}
\]
The analogous degree-36 triple representation is also stable.
These are necessary conditions for actual full GT. They are
not asserted for every actual gentle element.

\begin{proof}
We first fix the rightmost-strand convention, including the
inverse in the Hurwitz action. Let $M_{n,1}\subset B_{n+1}$
be the subgroup fixing the last strand. Forgetting that strand
gives the split Artin sequence
\[
 1\longrightarrow F_n\longrightarrow M_{n,1}
   \longrightarrow B_n\longrightarrow1,
 \tag{D23.P18.3}
\]
whose section uses $\sigma_1,\ldots,\sigma_{n-1}$. Its free
kernel has the peripheral basis
\[
 u_i=\sigma_n\cdots\sigma_{i+1}\sigma_i^2
          \sigma_{i+1}^{-1}\cdots\sigma_n^{-1},
 \qquad 1\le i\le n.
\]
In particular $u_n=\sigma_n^2$, and the kernel is its normal
closure in $M_{n,1}$. If $\alpha_\beta$ denotes conjugation by
the section element $\beta$ on $F_n$, the Artin relations give
\[
\begin{split}
 \alpha_{\sigma_i}(u_i)&=u_{i+1},\\
 \alpha_{\sigma_i}(u_{i+1})&=u_{i+1}^{-1}u_i u_{i+1},\\
 \alpha_{\sigma_i}(u_j)&=u_j\quad(j\ne i,i+1).
\end{split}                                                 \tag{D23.P18.4}
\]
Consequently the chosen left Hurwitz action on a homomorphism
$\varphi:F_n\to G$ is
\[
 q_n(\beta)[\varphi]=[\varphi\circ\alpha_\beta^{-1}].
 \tag{D23.P18.5}
\]
This is precisely the tuple action in the statement.

The sequence (D23.P18.3) remains a split exact sequence after
full profinite completion. Here is a direct topology check.
Characteristic finite-index subgroups of the finitely generated
free kernel are cofinal among its finite-index subgroups. For
such a subgroup $K$, the action on $F_n/K$ factors through a
finite quotient of $B_n$. The associated finite semidirect
product is a quotient of $M_{n,1}$ whose restriction to $F_n$
has kernel $K$. Thus the induced topology on $F_n$ is its full
profinite topology. The completed section and kernel have
trivial intersection, and their compact product contains the
dense discrete semidirect product; hence it is all of
$\widehat M_{n,1}$. This proves the asserted completed split
sequence. Since $M_{n,1}$ has finite index in $B_{n+1}$, its
completion is the corresponding open subgroup of
$\widehat B_{n+1}$.

We use the full-profinite Drinfeld braid automorphism theorem,
with its standard formula
\[
 \Phi_h(\sigma_1)=\sigma_1^\lambda,
 \qquad
 \Phi_h(\sigma_i)=f(y_i,\sigma_i^2)^{-1}
        \sigma_i^\lambda f(y_i,\sigma_i^2),
\quad
 y_i=\sigma_{i-1}\cdots\sigma_2\sigma_1^2
            \sigma_2\cdots\sigma_{i-1}.
 \tag{D23.P18.6}
\]
The action induces the identity on the symmetric-group
quotient, because $\lambda$ is odd and the conjugators are
pure. It preserves the section $\widehat B_n$ and restricts
there to its standard Drinfeld action. It therefore preserves
$\widehat M_{n,1}$.

Let $N=\widehat F_n$. Under the rightmost forgetting map,
the element $f(y_n,\sigma_n^2)$ maps to $f(y_n,1)=1$.
Thus this conjugator belongs to $N$, and
$\Phi_h(u_n)$ is $N$-conjugate to $u_n^\lambda$.
Since $N$ is the closed normal closure of $u_n$ in
$\widehat M_{n,1}$, $\Phi_h(N)\subseteq N$. Applying the same
argument to $h^{-1}$ gives equality. Put
$\kappa_h=\Phi_h|N$.

Every $u_i$ is a conjugate of $u_n$ by a braid in the section;
one may take $v_i=\sigma_i^{-1}\cdots\sigma_{n-1}^{-1}$ and
$u_i=v_i u_n v_i^{-1}$. The Artin action of any braid carries
a peripheral generator to an $N$-conjugate of the generator
with the corresponding puncture label. Since $\Phi_h(v_i)$
has the same permutation as $v_i$, it follows that
$\kappa_h(u_i)$ is $N$-conjugate to $u_i^\lambda$.

Define
\[
 V_h[\varphi]=[\varphi\circ\kappa_h^{-1}]
 \quad\text{on }\operatorname{Hom}_{\rm cont}(N,G)/G.
 \tag{D23.P18.7}
\]
This is a bijection, and when $\lambda=1$ it preserves the
condition $\varphi(u_i)\in\mathcal C$ for all $i$. Thus it
restricts to a permutation of the entire finite set
$\mathcal C^n/G$.
The semidirect product gives the exact covariance
\[
 \kappa_h\alpha_\beta\kappa_h^{-1}
       =\alpha_{\Phi_h(\beta)}.
\]
Combining it with (D23.P18.5)--(D23.P18.7) proves (D23.P18.1)
with the displayed orientation. If a different section is
used and a kernel cocycle appears, the covariance differs
only by an inner automorphism of $N$. Precomposition by that
inner automorphism becomes simultaneous conjugation by its
$G$-image, and hence disappears on $\operatorname{Hom}(N,G)/G$.

Because $\Phi_h$ is onto $\widehat B_n$, equation
(D23.P18.1) carries each $B_n$-orbit bijectively onto another
$B_n$-orbit of the same degree. A degree occurring exactly
once must therefore be preserved. Restriction of (D23.P18.1)
to it gives permutation conjugacy and kernel invariance.

For D23.P17 the complete four-tuple orbit degree list is
\[
 1,4,9,9,9,12,12,12,12,72,324,360,405,
 1344,1512,1512,2187.
\]
The degree 1344 therefore occurs exactly once in the entire
$\mathcal C^4/G$. The complete triple list is
\[
 1,1,2,2,4,4,4,4,4,36,40,60,60,
\]
so the degree 36 is likewise isolated.

Finally evaluate (D23.P18.6) on the three generators of
$B_4$. For $i=2$, $y_2=\sigma_1^2$, giving $F^{-1}BF$.
For $i=3$, replace $y_3$ by $z_3=(\sigma_1\sigma_2)^3$:
indeed $z_3=\sigma_1^2y_3$, and $\sigma_1^2$ commutes with
both $y_3$ and $\sigma_3^2$. A commutator word has zero
exponent sum in its first input, so this commuting factor
does not change its value. Thus the conjugator is
$f(X_4,X_3)=f(X_3,X_4)^{-1}=R^{-1}$, by the two-cycle.
This gives exactly (D23.P18.2).
\end{proof}

\paragraph{Additional functorial restrictions.}
The permutation (D23.P18.7) commutes with postcomposition by
every automorphism of $G$ preserving $\mathcal C$. It preserves
the image subgroup of a homomorphism because $\kappa_h$ is
onto. It also preserves the conjugacy class of the peripheral
product when $\lambda=1$: the boundary word $u_1\cdots u_n$
is $z_n^{-1}z_{n+1}$, and the standard action sends each full
twist $z_j$ to $z_j^\lambda$. These restrictions may be imposed
in a simultaneous-conjugacy test; no extra geometric rigidity
of $G$ or of a matrix representation is being assumed.

\paragraph{Primary source and check scope.}
The full-profinite automorphism theorem used in
(D23.P18.6) is Lochak--Schneps, \emph{The
Grothendieck--Teichm\"uller group and automorphisms of braid
groups}, LMS Lecture Notes 200 (1994), \S5, Proposition 10,
author-PDF p.23 (PDF index 22):
\url{https://webusers.imj-prg.fr/~leila.schneps/Fls.pdf}.
Its formula uses the displayed $y_i$ marking. Nakamura,
RIMS-1691, \S5.4, equation (5.4.2), author-PDF p.28,
\url{https://www.kurims.kyoto-u.ac.jp/preprint/file/RIMS1691.pdf},
records the equivalent partial-full-twist formula in Galois
scope; that source is used to check the marking, not to
replace the full-GT theorem by a Galois statement.
The new proof here is the induced finite-Hurwitz stability
argument and its application to the two completely enumerated
isolated orbits.

The independent exact word replay
\path{hurwitz_stability_artin_check.py} verifies
(D23.P18.4), all displayed peripheral transports, the boundary
full-twist relation, and the commuting input replacement in
$B_4$ and $B_5$, using the faithful Artin representation.
Its JSON keeps the literal peripheral braid words. The
published full-GT automorphism theorem is a cited input, not
independently reproved by that finite word check.
